\section{Method}

\begin{table*}[t]
    \caption{Quantitative Results. Performance comparison of fine-tuned and zero-shot models across key metrics.}
    \label{tab:metrics}
    \centering
    \begin{tabular}{p{4cm} wc{1.2cm} wc{1.2cm} wc{1.2cm} wc{1.2cm} wc{1.2cm} wc{1.2cm} wc{2cm}}
        \toprule
         & \thead{Pattern \\ Match}$\uparrow$ & \thead{JSON \\ Validity}$\uparrow$ & \thead{ID \\ IoU}$\uparrow$ & \thead{Floor \\ Error}$\downarrow$ & \thead{Area \\ Error}$\downarrow$ & \thead{Site \\ IoU}$\uparrow$ & \thead{Contextual \\ Relevance}$\uparrow$ \\
        \midrule
        CoMa-2B & 0.84 & 0.63 & 0.46 & 0.76 & 1.30 & 0.01 & 0.15 \\
        CoMa-4B & 0.95 & 0.72 & 0.71 & 0.49 & 2.86 & 0.03 & 0.18 \\
        CoMa-8B & 0.94 & 0.79 & 0.75 & 0.42 & 1.90 & 0.05 & 0.24 \\
        \midrule
        Qwen3-VL-235B-Instruct & \textbf{1.00} & \textbf{0.99} & \textbf{0.99} & \textbf{0.12} & \textbf{0.79} & \textbf{0.10} & \textbf{0.25} \\
        \bottomrule
    \end{tabular}
  \end{table*}

We frame the problem of building massing generation as a conditional vision-language task. The objective is to generate a 3D massing model conditioned on a textual description of its functional-economical requirements, development site contour, and its urban visual context.

\subsection{Finetuned VLM Approach}
We fine-tune the Qwen3-VL series \cite{Qwen-VL} of Vision Language Models on our CoMa-20K dataset. The model input is a concatenated prompt consisting of the site's \texttt{requirements} formatted as a text string, the \texttt{site\_contour} coordinates, and a contextual image (\texttt{env\_render\_high}). The training objective is to autoregressively generate the tokenized 3D massing geometry. We experiment with three model variants (2B, 4B, and 8B parameters) to study the scalability of this approach.

\subsection{Zero-shot VLM Approach}
To benchmark the inherent capability of large VLMs on this task, we employ \textit{Qwen3-VL-235B-A22B-Instruct} \cite{huggingface} in a zero-shot setting. The model is prompted with detailed architectural rules and spatial logic constraints, including:

\begin{itemize}
    \item \textbf{Spatial Logic:} Strict adherence to site boundaries, coordinate system consistency, and geometric containment
    \item \textbf{Placement Logic:} Footprint area estimation, orientation parallel to site contours, and minimum buffer zones
    \item \textbf{Geometric Articulation:} Avoidance of simple rectangles through L/U-shaped plans, step-backs, and varied roof heights
    \item \textbf{Structural Constraints:} Proper elevation calculations based on floor counts and heights
\end{itemize}

The prompt explicitly defines the input-output format and emphasizes outputting only structured JSON geometry, enabling direct comparison with our fine-tuned approaches while leveraging the model's pre-trained capabilities.

